\subsection{Estructuras de Datos: \texttt{class} frente a \texttt{record}}

\subsubsection{Regla: mutabilidad como criterio de elección}

\textbf{Regla:} Entidades cuyo estado cambia durante la partida se declaran \texttt{class}. Valores inmutables o eventos que solo describen un hecho ocurrido se declaran \texttt{record} (Microsoft, 2026b).

\textbf{Código correcto:}
\textit{Archivo \texttt{Player.cs}:}
\begin{lstlisting}[style=csharpstyle]
public sealed class Player
{
    public Player(string name)
    {
        Name = name;
        Score = 0;
    }

    public string Name { get; }
    public int Score { get; private set; }

    public void AddScore(int points)
    {
        Score += points;
    }
}
\end{lstlisting}

\textit{Archivo \texttt{DiceRolled.cs}:}
\begin{lstlisting}[style=csharpstyle]
public sealed record DiceRolled(Player Player, int FaceValue);
\end{lstlisting}

\textbf{Código incorrecto:}
\textit{Archivo \texttt{Player.cs}:}
\begin{lstlisting}[style=csharpstyle]
public sealed record Player
{
    public Player(string name)
    {
        Name = name;
        Score = 0;
    }

    public string Name { get; }
    public int Score { get; private set; }

    public void AddScore(int points)
    {
        Score += points;
    }
}
\end{lstlisting}

\textit{Archivo \texttt{DiceRolled.cs}:}
\begin{lstlisting}[style=csharpstyle]
public sealed class DiceRolled
{
    public DiceRolled(Player player, int faceValue)
    {
        Player = player;
        FaceValue = faceValue;
    }

    public Player Player { get; }
    public int FaceValue { get; }
}
\end{lstlisting}
